\documentclass[12pt]{tietreport}

% =========================================================
% PACKAGES
% =========================================================

\usepackage{graphicx}
\usepackage{amsmath}
\usepackage{amssymb}
\usepackage{array}
\usepackage{booktabs}
\usepackage{longtable}
\usepackage{float}
\usepackage{enumitem}
\usepackage{xurl}

% =========================================================
% BIBLIOGRAPHY
% =========================================================

\usepackage[
    date=year,
    style=alphabetic,
    backend=biber,
    sorting=nty,
    maxnames=5,
    minnames=3,
    maxalphanames=4,
    minalphanames=3,
    backref=true,
    doi=false,
    isbn=false,
    url=false,
    eprint=false
]{biblatex}

\addbibresource{references.bib}

% =========================================================
% DOCUMENT INFORMATION
% =========================================================

\TitlePageHeader{UCS 503 - Software Engineering}

\ProjectTitle{ICU Patient Face Monitoring System}

\ProjectSubTitle{
A Computer-Vision-Based Non-Contact System for Continuous
Facial and Ocular Monitoring of Intensive Care Unit Patients
}

\author{
    \texttt{1024240103} Pahul Singh\\
    \texttt{1024240039} Harjot Singh
}

\TitlePageSubText{
    BE Third Year, Data Science and Artificial Intelligence
}

\AdvisorName{
    Ms. Anushka, UCS 503
}

\TitlePageFooterText{
    \large Project Proposal \\[0.2cm]
    \bfseries Computer Science and Engineering Department\\[0.15cm]
    Thapar Institute of Engineering and Technology, Patiala
}

% =========================================================
% DOCUMENT
% =========================================================

\begin{document}

\maketitle

% =========================================================
% TABLE OF CONTENTS
% =========================================================

\tableofcontents

\newpage

% =========================================================
% CHAPTER 1
% =========================================================

\chapter{Introduction}

\section{Background and Context}

Intensive Care Units (ICUs) contain patients who require continuous
observation because their physical and neurological condition can
change rapidly. Traditional patient monitoring systems primarily
depend on sensors attached to the patient, such as ECG electrodes,
pulse oximeters, blood-pressure cuffs, and other clinical devices.
Although these systems provide important physiological measurements,
they do not directly describe visible facial and ocular behaviour.

Computer vision provides a non-contact method for observing visible
patient characteristics through a camera. A camera can capture facial
information without requiring additional physical sensors to be
attached to the patient's face. Facial landmarks can then be used to
locate important regions such as the eyes and determine geometric
features associated with eye opening, closing, blinking, and eye
movement.

Vision-based patient monitoring has been studied as a possible
supplement to conventional monitoring systems
\cite{bevilacqua2016visionreview,davoudi2019intelligenticu}.
The purpose of this project is to develop a focused Face + Eye
monitoring pipeline that converts camera frames into structured
information that can support continuous ICU observation.

The system is designed as a supplementary monitoring and
decision-support system. It is not intended to replace medical
professionals or provide an independent medical diagnosis.

\section{Problem Statement}

Continuous visual observation of an ICU patient can be difficult when
healthcare staff must simultaneously supervise multiple patients.
Important changes in facial and ocular behaviour may not always be
immediately noticed.

The project therefore addresses the following problem:

\begin{quote}
How can facial and eye information obtained from a standard camera
be automatically processed into meaningful, time-based monitoring
values and alerts for an ICU patient without requiring physical
contact with the patient's face?
\end{quote}

The proposed system focuses specifically on the face and eye regions.
It detects the patient's face, identifies facial and ocular landmarks,
extracts geometric features, analyses their behaviour over time, and
converts the resulting information into global monitoring values and
meaningful alerts.

\section{Project Objectives}

The main objectives of the project are:

\begin{enumerate}[leftmargin=*]

    \item Detect a patient's face from a live or recorded camera
    stream.

    \item Locate facial landmarks required for accurate eye-region
    analysis.

    \item Extract landmarks corresponding to the left and right eyes.

    \item Calculate geometric eye features such as the Eye Aspect
    Ratio (EAR).

    \item Identify changes in eye state over time.

    \item Analyse eye closure and blinking behaviour using temporal
    information.

    \item Extract useful information related to eye movement and
    approximate gaze behaviour where reliable landmarks are
    available.

    \item Convert frame-level measurements into stable temporal
    values.

    \item Generate global monitoring values representing the current
    state of the facial and ocular observations.

    \item Generate meaningful alerts when predefined abnormal or
    noteworthy patterns are detected.

    \item Provide processed values to the monitoring interface so that
    the information can be displayed in an understandable form.

\end{enumerate}

% =========================================================
% CHAPTER 2
% =========================================================

\chapter{Methodology and Design}

\section{System Architecture}

\subsection{High-Level Design}

The proposed system follows a sequential computer-vision pipeline.
The major stages are:

\begin{center}
\begin{verbatim}
Camera / Video Input
        |
        v
Video Frame Acquisition
        |
        v
Face Detection
        |
        v
Facial Landmark Detection
        |
        v
Eye Landmark Extraction
        |
        v
Eye Feature Extraction
        |
        v
Temporal Processing
        |
        +-------------------+
        |                   |
        v                   v
 Global Values          Alert Generation
        |                   |
        +---------+---------+
                  |
                  v
        Monitoring Interface
\end{verbatim}
\end{center}

The camera provides a continuous sequence of frames. Each frame is
processed to locate the patient's face and facial landmarks. The eye
landmarks are then used to calculate ocular features. These
frame-level measurements are not directly treated as final outputs.
Instead, they are processed over time to reduce noise and identify
meaningful patterns.

The processed information is converted into global values and alert
states. These results are then supplied to the monitoring interface,
where they can be presented to a user in an understandable format.

\section{Technical Stack}

The proposed implementation uses the following technologies:

\begin{table}[H]
\centering
\caption{Proposed Technical Stack}
\begin{tabular}{p{0.28\textwidth}p{0.62\textwidth}}
\toprule
\textbf{Technology} & \textbf{Purpose} \\
\midrule
Python & Main programming language for the computer-vision pipeline \\
OpenCV & Video capture, frame processing and image operations \\
MediaPipe Face Mesh & Facial landmark detection and face geometry \\
NumPy & Numerical calculations and feature processing \\
Flask/FastAPI & Optional service layer for delivering processed values \\
Browser Interface & Display of monitoring values and alerts \\
Git & Version control and project collaboration \\
\bottomrule
\end{tabular}
\end{table}

OpenCV provides a widely used computer-vision framework for image and
video processing. MediaPipe Face Mesh provides real-time facial
landmark information suitable for extracting geometric information
from the face \cite{kartynnik2019facemesh,opencv,mediapipe}.

\section{Implementation Methodology}

\subsection{Face Detection}

The first stage of the pipeline identifies the patient's face in
each incoming frame.

The face detection stage produces a region of interest containing
the face. Processing can then be restricted primarily to this region,
reducing unnecessary computation from the background.

Traditional face detection methods such as the Viola--Jones
framework demonstrated the possibility of real-time face detection
using visual features and classifiers \cite{viola2004rapid}.
Modern landmark-based approaches can provide additional geometric
information required by the remaining stages.

\subsection{Facial Landmark Detection}

After locating the face, facial landmarks are detected.

A facial landmark is a point corresponding to a meaningful position
on the face. Multiple landmarks collectively describe the geometry of
the face. The eye regions can be extracted from these landmarks.

The landmark stage provides a common coordinate representation that
can be used for further measurements instead of relying only on raw
pixel appearance.

\subsection{Eye Landmark Extraction}

The left and right eye regions are extracted from the complete set of
facial landmarks.

For each eye, the relevant landmarks describe the upper eyelid,
lower eyelid, and horizontal eye boundaries. These points provide the
geometric information required for estimating whether the eye is
open, partially closed, or closed.

\subsection{Eye Feature Extraction}

One of the primary geometric features is the Eye Aspect Ratio (EAR).
EAR provides a numerical representation of the vertical opening of an
eye relative to its horizontal width.

For six eye landmarks, EAR can be represented as:

\begin{equation}
EAR =
\frac{
\left\|p_2-p_6\right\|
+
\left\|p_3-p_5\right\|
}{
2\left\|p_1-p_4\right\|
}
\end{equation}

where $p_1,\ldots,p_6$ represent selected eye landmarks and
$\|\cdot\|$ represents Euclidean distance.

When an eye is open, the vertical distances are generally larger.
When the eye closes, these distances decrease. This makes EAR useful
for analysing eye closure and blinking behaviour
\cite{soukupova2016eyeblink}.

\subsection{Eye-State Analysis}

The calculated eye features can be converted into an eye-state
classification.

A simplified representation is:

\begin{center}
\begin{tabular}{c|c}
\textbf{Condition} & \textbf{Possible Eye State}\\
\hline
EAR above threshold & Open\\
EAR near threshold & Partially closed\\
EAR below threshold & Closed
\end{tabular}
\end{center}

The threshold should not be considered a universal medical value.
It may depend on camera position, patient characteristics, lighting,
facial orientation, and the specific implementation.

\subsection{Blink Analysis}

A blink can be interpreted as a short-duration transition from an
open-eye state to a closed-eye state and then back to an open-eye
state.

The system can maintain a temporal sequence of eye states:

\begin{verbatim}
Open -> Closing -> Closed -> Opening -> Open
\end{verbatim}

This sequence can be used to identify individual blink events.

The duration and frequency of such events can be calculated from the
sequence of observations. Real-time landmark-based blink detection has
previously been investigated using facial landmark geometry
\cite{soukupova2016eyeblink}.

\subsection{Eye Movement and Gaze Analysis}

Facial and eye landmarks can also provide information about relative
eye movement.

Instead of attempting to determine an exact medical gaze direction,
the system can represent approximate movement categories such as:

\begin{itemize}[leftmargin=*]
    \item left,
    \item right,
    \item upward,
    \item downward,
    \item approximately centred,
    \item uncertain.
\end{itemize}

The reliability of this information depends on camera placement,
head orientation, illumination, eyelid visibility, and landmark
quality.

\subsection{Temporal Processing}

Frame-by-frame measurements can be noisy. A single incorrect
landmark position should therefore not immediately generate an alert.

The system maintains a short temporal history of observations.

For a sequence:

\[
x_1,x_2,\ldots,x_n
\]

a simple moving average can be calculated as:

\begin{equation}
\bar{x}_t =
\frac{1}{N}
\sum_{i=0}^{N-1}x_{t-i}
\end{equation}

where $N$ is the number of recent observations.

Temporal processing can therefore make the output more stable and
reduce false alerts caused by individual noisy frames.

\subsection{Global Value Generation}

The system converts multiple low-level measurements into a smaller
set of global monitoring values.

Examples include:

\begin{itemize}[leftmargin=*]
    \item current eye state,
    \item left-eye EAR,
    \item right-eye EAR,
    \item average EAR,
    \item blink count,
    \item recent blink frequency,
    \item eye-closure duration,
    \item approximate gaze state,
    \item confidence or validity state,
    \item current alert state.
\end{itemize}

These values provide a compact representation of the patient's
recent facial and ocular observations.

\subsection{Alert Generation}

Alerts are generated from processed temporal values rather than
individual raw frames.

For example, if the eye remains below the closure threshold for a
predefined period, the system can generate a prolonged-eye-closure
alert.

A conceptual rule is:

\begin{equation}
Alert =
\begin{cases}
1, & \text{if closure duration} > T_c\\
0, & \text{otherwise}
\end{cases}
\end{equation}

where $T_c$ represents a configurable duration threshold.

Similarly, abnormal or noteworthy changes in blink behaviour can be
represented by additional rules.

These alerts are intended to act as supplementary monitoring
indicators rather than clinical diagnoses.

% =========================================================
% CHAPTER 3
% =========================================================

\chapter{System Requirements}

\section{Hardware Requirements}

The proposed system requires:

\begin{itemize}[leftmargin=*]
    \item A standard RGB camera or webcam.
    \item A computer or workstation capable of real-time video
    processing.
    \item Sufficient RAM for the operating system and computer-vision
    pipeline.
    \item Optional GPU acceleration for higher-resolution or
    multi-stream processing.
    \item Stable camera positioning near the patient's bed.
\end{itemize}

\section{Software Requirements}

\begin{table}[H]
\centering
\caption{Software Requirements}
\begin{tabular}{p{0.30\textwidth}p{0.58\textwidth}}
\toprule
\textbf{Component} & \textbf{Requirement} \\
\midrule
Operating System & Windows/Linux/macOS or equivalent environment \\
Programming Language & Python \\
Computer Vision & OpenCV \\
Landmark Detection & MediaPipe Face Mesh \\
Numerical Processing & NumPy \\
Backend & Flask or FastAPI, if required \\
Interface & Browser-based monitoring interface \\
Version Control & Git \\
\bottomrule
\end{tabular}
\end{table}

\section{Functional Requirements}

The system shall:

\begin{enumerate}[leftmargin=*]
    \item Accept live or recorded video input.
    \item Detect the patient's face.
    \item Detect facial landmarks.
    \item Identify eye landmarks.
    \item Calculate eye-related geometric features.
    \item Estimate eye state.
    \item Detect blink events.
    \item Maintain temporal information.
    \item Calculate global monitoring values.
    \item Generate configurable alerts.
    \item Make processed values available to the monitoring interface.
\end{enumerate}

\section{Non-Functional Requirements}

\subsection{Real-Time Performance}

The pipeline should process frames with sufficiently low latency for
continuous monitoring. The target is to maintain approximately
15 frames per second or higher where the available hardware permits.

\subsection{Reliability}

The system should tolerate temporary landmark loss, lighting
variation, and small head movements without crashing or generating
continuous false alerts.

\subsection{Accuracy}

The system should be evaluated using labelled video sequences and
appropriate metrics rather than assuming that every detected
condition is correct.

\subsection{Maintainability}

The face detection, eye analysis, temporal processing, alert
generation, and interface communication components should remain
modular so that individual components can be improved independently.

\subsection{Privacy}

Patient video represents sensitive information. The system should
prefer local processing where possible and should avoid unnecessary
storage or transmission of raw video.

% =========================================================
% CHAPTER 4
% =========================================================

\chapter{Data Processing and Alert System}

\section{Data Flow}

The complete data flow can be represented as:

\begin{center}
\begin{verbatim}
RAW VIDEO
   |
   v
Frame Acquisition
   |
   v
Face Detection
   |
   v
Facial Landmarks
   |
   v
Eye Landmarks
   |
   v
Geometric Features
   |
   v
Temporal Processing
   |
   +--------------------+
   |                    |
   v                    v
Global Values       Alert Engine
   |                    |
   +---------+----------+
             |
             v
      Monitoring Interface
\end{verbatim}
\end{center}

\section{Raw Data}

Raw input consists primarily of camera frames.

Each frame contains pixel information representing the patient's
visible face and surrounding environment.

Raw frames may also contain irrelevant information such as:

\begin{itemize}[leftmargin=*]
    \item background objects,
    \item bedding,
    \item medical equipment,
    \item other visual noise.
\end{itemize}

The processing pipeline therefore extracts only the information
needed for face and eye monitoring.

\section{Extracted Data}

The face and landmark stages produce structured information such as:

\begin{itemize}[leftmargin=*]
    \item face location,
    \item facial landmark coordinates,
    \item eye landmark coordinates,
    \item eye-region geometry,
    \item detection confidence or validity information.
\end{itemize}

This information is significantly smaller and more structured than
the original video frame.

\section{Processed Data}

The extracted measurements are processed to generate:

\begin{itemize}[leftmargin=*]
    \item left-eye EAR,
    \item right-eye EAR,
    \item average eye opening,
    \item eye state,
    \item blink events,
    \item closure duration,
    \item temporal statistics,
    \item approximate eye movement,
    \item confidence information.
\end{itemize}

Temporal processing is important because an individual frame may
contain an incorrect or incomplete observation.

\section{Global Values}

Global values summarize the current condition of the Face + Eye
pipeline.

A possible global monitoring record can be represented as:

\begin{verbatim}
{
    eye_state,
    left_ear,
    right_ear,
    average_ear,
    blink_count,
    closure_duration,
    gaze_state,
    confidence,
    alert_state
}
\end{verbatim}

The exact set of fields can be changed according to the final
implementation.

\section{Meaningful Alerts}

The purpose of the alert system is to convert numerical measurements
into understandable monitoring information.

Possible alert categories include:

\begin{table}[H]
\centering
\caption{Example Alert Categories}
\begin{tabular}{p{0.30\textwidth}p{0.55\textwidth}}
\toprule
\textbf{Alert} & \textbf{Meaning} \\
\midrule
Prolonged Eye Closure & Eye closure persists beyond a configurable duration \\
Repeated Closure & Frequent closure events are observed in a recent time window \\
Unusual Blink Pattern & Blink behaviour differs from a configured baseline or rule \\
Low Confidence & Face or eye landmarks are unreliable \\
Face Not Detected & The patient's face cannot currently be located \\
\bottomrule
\end{tabular}
\end{table}

Alerts should contain both the event and sufficient context for a
user to understand why the alert was produced.

% =========================================================
% CHAPTER 5
% =========================================================

\chapter{Expected Results}

\section{Expected Processing Output}

The expected output of the Face + Eye pipeline is a continuously
updated collection of structured monitoring values.

For every processed observation, the system should be able to provide
information similar to:

\begin{table}[H]
\centering
\caption{Expected Monitoring Output}
\begin{tabular}{p{0.32\textwidth}p{0.50\textwidth}}
\toprule
\textbf{Output} & \textbf{Description} \\
\midrule
Face Status & Whether a usable face is currently detected \\
Eye State & Open, partially closed, closed, or uncertain \\
EAR & Numerical measure of eye opening \\
Blink Information & Detection/count of blink events \\
Closure Duration & Duration of continuous or recent closure \\
Gaze State & Approximate relative eye movement state \\
Confidence & Reliability of the current observation \\
Alert & Current meaningful monitoring alert \\
\bottomrule
\end{tabular}
\end{table}

The final processed values are intended to be sent to the monitoring
interface. The interface can then convert these values into
readable indicators, status cards, charts, or alerts.

% =========================================================
% CHAPTER 6
% =========================================================

\chapter{Applications}

The proposed Face + Eye monitoring system may be useful in:

\begin{enumerate}[leftmargin=*]

    \item \textbf{ICU Monitoring:}
    Supplementary observation of facial and ocular behaviour.

    \item \textbf{Continuous Observation:}
    Automated monitoring can provide a persistent stream of
    computer-generated observations.

    \item \textbf{Remote Monitoring:}
    Processed values may be made available to an authorised
    monitoring interface.

    \item \textbf{Research:}
    The system can provide structured facial and eye measurements
    for research involving patient behaviour.

    \item \textbf{Alert Support:}
    The system can highlight potentially noteworthy changes for
    further human attention.

\end{enumerate}

The system should always be treated as a supplementary tool and not
as a replacement for professional medical judgement.

% =========================================================
% CHAPTER 7
% =========================================================

\chapter{Advantages and Limitations}

\section{Advantages}

\begin{itemize}[leftmargin=*]

    \item \textbf{Non-contact:}
    No physical sensor is required on the patient's face.

    \item \textbf{Continuous:}
    Video can be processed continuously rather than through
    occasional manual observations.

    \item \textbf{Automated:}
    Facial and ocular measurements can be calculated automatically.

    \item \textbf{Structured Output:}
    Raw visual information is converted into numerical values and
    understandable alerts.

    \item \textbf{Modular:}
    Individual processing stages can be independently improved.

    \item \textbf{Interface Integration:}
    Processed global values can be supplied to a monitoring
    interface.

\end{itemize}

\section{Limitations}

The system has several limitations:

\begin{itemize}[leftmargin=*]

    \item Poor illumination can reduce landmark accuracy.

    \item Significant head rotation can make eye landmarks unreliable.

    \item Occlusion caused by blankets, equipment, hands, or other
    objects can interfere with detection.

    \item Camera placement affects the quality of measurements.

    \item Individual differences in facial structure can affect
    threshold-based eye-state classification.

    \item A computer-vision alert does not necessarily indicate a
    medical emergency.

    \item Raw video may contain sensitive patient information and
    therefore requires appropriate privacy controls.

\end{itemize}

% =========================================================
% CHAPTER 8
% =========================================================

\chapter{Testing and Evaluation}

\section{Testing Methodology}

The system should be evaluated component by component before
performing complete end-to-end testing.

\subsection{Face Detection Testing}

Face detection should be tested under:

\begin{itemize}[leftmargin=*]
    \item different lighting conditions,
    \item different face orientations,
    \item different distances from the camera,
    \item partial occlusion,
    \item periods of temporary detection failure.
\end{itemize}

\subsection{Eye Detection Testing}

Eye landmarks should be evaluated for:

\begin{itemize}[leftmargin=*]
    \item open eyes,
    \item closed eyes,
    \item blinking,
    \item partially visible eyes,
    \item head movement,
    \item varying illumination.
\end{itemize}

\subsection{Feature Testing}

EAR calculations should be checked against manually labelled eye
states.

The calculated values should behave consistently with the expected
geometric changes of the eye.

\subsection{Temporal Testing}

Temporal processing should be tested using sequences containing:

\begin{itemize}[leftmargin=*]
    \item normal eye opening,
    \item individual blinks,
    \item multiple blinks,
    \item prolonged closure,
    \item temporary landmark loss.
\end{itemize}

\subsection{Alert Testing}

Alert generation should be evaluated using controlled test
sequences.

The test should determine whether:

\begin{enumerate}[leftmargin=*]
    \item valid events generate alerts,
    \item normal observations avoid unnecessary alerts,
    \item temporary detection errors do not create excessive alerts,
    \item alert states return to normal when the condition ends.
\end{enumerate}

\subsection{End-to-End Testing}

Finally, the complete system should be tested from camera input to
monitoring interface.

The complete path is:

\begin{center}
\begin{verbatim}
Camera
  -> Frame
  -> Face
  -> Landmarks
  -> Eye Features
  -> Temporal Processing
  -> Global Values
  -> Alerts
  -> Monitoring Interface
\end{verbatim}
\end{center}

\section{Evaluation Metrics}

The following metrics can be used during evaluation:

\begin{table}[H]
\centering
\caption{Proposed Evaluation Metrics}
\begin{tabular}{p{0.30\textwidth}p{0.55\textwidth}}
\toprule
\textbf{Metric} & \textbf{Purpose} \\
\midrule
Detection Accuracy & Measures correct face/eye detections \\
Precision & Measures proportion of predicted events that are correct \\
Recall & Measures proportion of actual events that are detected \\
F1 Score & Balances precision and recall \\
Latency & Measures processing delay \\
FPS & Measures real-time processing throughput \\
False Alert Rate & Measures unnecessary alerts \\
Miss Rate & Measures important events that were not detected \\
\bottomrule
\end{tabular}
\end{table}

The project may use targets such as:

\begin{itemize}[leftmargin=*]
    \item latency below approximately 50 ms per processing step where
    hardware permits,
    \item face detection accuracy above 95\% on the selected test
    data,
    \item blink detection accuracy above 90\% on labelled sequences,
    \item processing throughput of at least approximately 15 FPS.
\end{itemize}

These values are proposed evaluation targets and should not be
presented as achieved performance until experimental testing has
confirmed them.

% =========================================================
% CHAPTER 9
% =========================================================

\chapter{Project Timeline}

\begin{table}[H]
\centering
\caption{Proposed Project Timeline}
\begin{tabular}{p{0.20\textwidth}p{0.62\textwidth}}
\toprule
\textbf{Phase} & \textbf{Work} \\
\midrule
Phase 1 & Requirement analysis and literature review \\
Phase 2 & Camera input and face detection \\
Phase 3 & Facial and eye landmark extraction \\
Phase 4 & Eye feature calculation and state analysis \\
Phase 5 & Temporal processing and blink analysis \\
Phase 6 & Global value and alert generation \\
Phase 7 & Monitoring interface integration \\
Phase 8 & Testing, evaluation and documentation \\
\bottomrule
\end{tabular}
\end{table}

The phases can overlap where appropriate. For example, interface
development can proceed while the computer-vision pipeline is being
tested.

% =========================================================
% CHAPTER 10
% =========================================================

\chapter{Future Scope}

Future development can extend the system in several directions.

\begin{enumerate}[leftmargin=*]

    \item \textbf{Improved Personalisation:}
    Patient-specific baselines could be used instead of relying only
    on fixed thresholds.

    \item \textbf{Advanced Temporal Models:}
    Machine-learning models could be used to identify more complex
    temporal patterns.

    \item \textbf{Multi-Camera Monitoring:}
    Multiple camera views could improve robustness against
    occlusion and head rotation.

    \item \textbf{Improved Gaze Estimation:}
    More sophisticated methods could provide more reliable relative
    gaze information.

    \item \textbf{Better Uncertainty Handling:}
    Confidence-aware processing could prevent unreliable frames from
    influencing alerts.

    \item \textbf{Historical Analysis:}
    Processed global values could be stored securely to allow
    authorised users to examine trends over time.

    \item \textbf{Scalable Deployment:}
    The system could be extended to support multiple patient
    monitoring streams.

\end{enumerate}

% =========================================================
% CHAPTER 11
% =========================================================

\chapter{Conclusion}

This project proposes a computer-vision-based Face + Eye monitoring
system for continuous, non-contact observation of ICU patients.

The proposed pipeline begins with camera input and proceeds through
face detection, facial landmark detection, eye landmark extraction,
eye feature calculation, and temporal processing. The resulting
information is converted into global monitoring values and meaningful
alerts.

The final processed information can then be provided to a monitoring
interface so that complex computer-vision measurements are presented
in a simple and understandable form.

The system is intended to supplement existing patient monitoring
rather than replace healthcare professionals. Its usefulness will
ultimately depend on the quality of the camera input, robustness of
landmark detection, accuracy of temporal analysis, appropriate
threshold selection, and systematic experimental evaluation.

% =========================================================
% REFERENCES
% =========================================================

\printbibliography

\end{document}